% ============================================================
%  Section 4.3 – Intervalles de confiance
% ============================================================

\subsection{Intervalles de confiance}

On construit trois intervalles de confiance au niveau $1-\alpha = 95\,\%$
($\alpha = 5\,\%$).
Les conditions de validité sont vérifiées pour chacun avant d'appliquer la formule.

% ── IC π1 ───────────────────────────────────────────────────────────────────
\paragraph{IC de la proportion $\pi_1$ -- sommeil insuffisant avant un DE (V6).}

\textit{Paramètre d'intérêt} : $\pi_1 = P(\text{étudiant·e dort} < 7\,\text{h avant un DE})$
(modalités 1, 2, 3 de V6).

\textit{Conditions} (Déf.~2.11) : $n = 68 \geq 30$ {\checkmark},
$n\hat{p}_1 = 55 \geq 5$ {\checkmark}, $n(1-\hat{p}_1) = 13 \geq 5$ {\checkmark}.

\[
  \hat{p}_1 = \frac{55}{68} \approx 0{,}8088
  \qquad
  \text{IC}_{95\%}(\pi_1)
  = \hat{p}_1 \pm z_{0{,}025}\sqrt{\frac{\hat{p}_1(1-\hat{p}_1)}{n}}
  = 0{,}8088 \pm 1{,}96 \times 0{,}0477
  \approx \boldsymbol{[0{,}7154\,;\,0{,}9023]}
\]

\textit{Interprétation} : on estime, avec un niveau de confiance de 95\,\%, qu'entre
71,5\,\% et 90,2\,\% des étudiant·es dorment moins de 7~heures la nuit précédant un DE.

% ── IC π2 ───────────────────────────────────────────────────────────────────
\paragraph{IC de la proportion $\pi_2$ -- assiduité totale aux CMs (V14).}

\textit{Paramètre d'intérêt} : $\pi_2 = P(\text{étudiant·e assiste à tous les CMs})$
(modalité « Tous les CMs » de V14).

\textit{Conditions} : $n = 68 \geq 30$ {\checkmark},
$n\hat{p}_2 = 29 \geq 5$ {\checkmark}, $n(1-\hat{p}_2) = 39 \geq 5$ {\checkmark}.

\[
  \hat{p}_2 = \frac{29}{68} \approx 0{,}4265
  \qquad
  \text{IC}_{95\%}(\pi_2)
  = 0{,}4265 \pm 1{,}96 \times 0{,}0600
  \approx \boldsymbol{[0{,}3089\,;\,0{,}5440]}
\]

\textit{Interprétation} : on estime qu'entre 30,9\,\% et 54,4\,\% des étudiant·es
assistent à tous leurs cours magistraux.
La valeur $1/3$ (référence du test T2) appartient à cet intervalle, ce qui
anticipe le résultat du test de conformité.

% ── IC m ────────────────────────────────────────────────────────────────────
\paragraph{IC de la moyenne $m$ -- annales travaillées avant CE/DE (V2).}

\textit{Paramètre d'intérêt} : $m = $ nombre moyen d'annales travaillées.

\textit{Conditions} (Déf.~2.14) : $n = 68 \geq 30$ {\checkmark} ;
$\sigma$ inconnue $\Rightarrow$ on utilise $S_e$ et la loi de Student
$\mathcal{T}(n-1) = \mathcal{T}(67)$ (§\,2.3.1).

\[
  \bar{x} \approx 2{,}1912,\quad
  S_e \approx 0{,}8853,\quad
  t_{67\,;\,0{,}025} \approx 1{,}9960
\]
\[
  \text{IC}_{95\%}(m)
  = \bar{x} \pm t_{67\,;\,0{,}025}\,\frac{S_e}{\sqrt{n}}
  = 2{,}1912 \pm 1{,}9960 \times 0{,}1074
  \approx \boldsymbol{[1{,}977\,;\,2{,}406]}
\]

\textit{Interprétation} : en moyenne, les étudiant·es travaillent entre
1,98 et 2,41 annales avant un examen, avec un niveau de confiance de 95\,\%.
La valeur $m_0 = 2$ appartient à cet intervalle : le test de conformité T3
ne devrait pas rejeter $H_0$.
